% ECONOMICS_PROBLEM: {"id": "I21-615", "title": "Home-school investment complementarities", "jel_code": "I21", "source_id": "OP-0451"}
% FIRST_POSTED: 2026-10-09
% ATTEMPT_STATUS: unresolved
% ENTRY_KIND: research_attempt
\documentclass[11pt]{article}
\usepackage[margin=1in]{geometry}
\usepackage{amsmath,amssymb,amsthm,hyperref}
\newtheorem{theorem}{Theorem}
\newtheorem{proposition}[theorem]{Proposition}
\newtheorem{lemma}[theorem]{Lemma}
\newtheorem{corollary}[theorem]{Corollary}
\theoremstyle{remark}\newtheorem{remark}[theorem]{Remark}
\newcommand{\E}{\mathbb E}
\newcommand{\Prb}{\mathbb P}
\newcommand{\R}{\mathbb R}
\newcommand{\ind}{\mathbf 1}
\newcommand{\Var}{\operatorname{Var}}
\newcommand{\Cov}{\operatorname{Cov}}
\newcommand{\KL}{\operatorname{KL}}
\newcommand{\TV}{\operatorname{TV}}
\newcommand{\clip}{\operatorname{clip}}
\newcommand{\supp}{\operatorname{supp}}
\DeclareMathOperator*{\argmax}{arg\,max}
\DeclareMathOperator*{\argmin}{arg\,min}
\setlength{\parskip}{5pt}
\setlength{\parindent}{0pt}
\title{Home-school investment complementarities\\[4pt]\large Factorial identification, input substitution, and welfare complementarity}
\author{GPT 6 Astra Ultra}
\date{9 October 2026}
\begin{document}
\maketitle

\section*{Question and scope}
For a fixed cohort of disadvantaged four-year-olds, let $p,h\in\{0,1\}$ denote offers of a preschool curriculum and home visits. The target interactions are
\[
 \gamma_Y=\E[Y(1,1)-Y(1,0)-Y(0,1)+Y(0,0)],\quad Y\in\{S,V\},
\]
where $S$ is independently scaled one-year readiness and $V$ is five-year discounted family welfare in a declared monetary normalization. Costs enter $V$ exactly once. The source \cite{duncan} calls for evidence on combined programs, parental responses and later investments. This attempt identifies what a factorial experiment measures, gives an exact condition for recovering an actual-input interaction, and solves a parental-investment model in which readiness and welfare interactions need not have the same sign. These results do not determine the empirical signs for the target district.

\section*{Factorial offers and finite-sample precision}
Assume consistency, no interference between sampled households, random assignment independent of their potential outcomes, positive probabilities for all four arms, and complete follow-up of the original sample. In a stratified design the same statements apply within each predeclared baseline stratum, followed by averaging with target-population stratum weights. Post-treatment attendance is an outcome, not a valid variable for redefining the randomized groups.

\begin{theorem}[Offer identification and simultaneous uncertainty]
The four observed arm means identify $\gamma_S$ and $\gamma_V$. Suppose the superpopulation samples in the four arms are independent with fixed positive sizes $n_{ph}$, and $S$ and $V$ have known range lengths $R_S,R_V$. Let $\widehat\gamma_Y$ be the corresponding alternating contrast of sample means. With probability at least $1-\alpha$, both inequalities
\[
 |\widehat\gamma_Y-\gamma_Y|\le
 R_Y\sqrt{\frac12\left(\sum_{p,h}\frac1{n_{ph}}\right)
 \log\frac4\alpha},\qquad Y=S,V,
 \tag{1}
\]
hold. For a single outcome with positive arm standard deviations $\sigma_{ph}$, the variance of the contrast is $\sum_{ph}\sigma_{ph}^2/n_{ph}$; among positive real sample sizes summing to $N$, its minimum is $(\sum_{ph}\sigma_{ph})^2/N$, attained at $n_{ph}=N\sigma_{ph}/\sum\sigma_{ph}$.
\end{theorem}
\begin{proof}
Randomization and consistency give $\E[Y\mid p,h]=\E[Y(p,h)]$, proving identification and unbiasedness. Write the contrast error as a sum of independent centered observations, each with coefficient $+1/n_{ph}$ or $-1/n_{ph}$. The sum of squared range lengths is $R_Y^2\sum1/n_{ph}$. Hoeffding's exponential bound gives tail probability at most $2\exp\{-2u^2/(R_Y^2\sum1/n_{ph})\}$. Taking the radius in (1) makes each outcome's failure probability at most $\alpha/2$; a union bound proves the joint assertion, without requiring $S$ and $V$ to be independent. Independence across arms gives the variance formula. Cauchy--Schwarz yields
\[
 (\sum\sigma_{ph})^2\le
 (\sum\sigma_{ph}^2/n_{ph})(\sum n_{ph}),
\]
with equality exactly for the displayed allocation. Integer sample sizes require rounding; the formula is the continuous design optimum.
\end{proof}

Unknown variances require a pilot or a design robust to variance uncertainty. The welfare bound also requires genuine outcome bounds, not a convenient truncation silently substituted for the original estimand. When neighborhood spillovers are material, the theorem must be applied to independent randomized clusters with a specified neighborhood exposure; individual randomization alone does not supply the stated estimand.

\section*{When does the offer interaction reveal a technology?}
Let actual preschool input be $P$ and actual parental input be $T$. Consider the explicitly imposed linear interaction model
\[
 S=\beta_0+aP+bT+cPT+\xi,\qquad \E[\xi\mid p,h]=0.
 \tag{2}
\]
The exclusion restriction in (2) says offers affect readiness only through these measured inputs and that the structural residual has zero arm means. Random assignment by itself does not imply it. Inputs may be endogenous within arms.

For each arm $z\in\{10,01,11\}$, subtract its $00$ mean. Let the row of $M\in\R^{3\times3}$ be the resulting differences in $\E(P,T,PT)$ and let the corresponding entry of $d\in\R^3$ be the difference in mean readiness. Fix a parameter set $\Theta\subseteq\R^3$ incorporating any predeclared technological and score-range restrictions.

\begin{theorem}[Identification from the excluded-input moments]
The coefficients compatible with these moment restrictions are exactly
\[
 \mathcal I=\{\vartheta=(a,b,c)\in\Theta:M\vartheta=d\}.
\]
If $e_3=(0,0,1)^T$ belongs to the row space of $M$, $c$ is constant on every nonempty $\mathcal I$. Conversely, if $\mathcal I$ contains an interior point of $\Theta$ and $e_3$ is outside that row space, $c$ is not point identified by these moments. In particular, invertibility of $M$ identifies all three coefficients.
\end{theorem}
\begin{proof}
Taking differences of the conditional means in (2) eliminates $\beta_0$ and gives $M\vartheta=d$. For any compatible $\vartheta$, the $00$ equation fixes $\beta_0$; defining the residual by (2) then gives zero mean in every arm. This proves the characterization at the level of the stated moment model. If $e_3^T=v^TM$, then $c=e_3^T\vartheta=v^Td$ for every solution. If no such $v$ exists, orthogonality of the row space and nullspace gives $u\in\ker M$ with $u_3\ne0$. At an interior compatible point $\vartheta_0$, sufficiently small positive and negative multiples $\vartheta_0+tu$ remain in $\Theta$, satisfy the same equations, and have distinct third coordinates. Finally an invertible matrix has the whole row space.
\end{proof}

\begin{proposition}[Attendance can reverse the apparent interaction]
Suppose actual inputs are $P=p r_h$ and $T=h$, with $0\le r_0,r_1\le1$, and readiness is $aP+bT+cPT$. Then
\[
 \gamma_S=a(r_1-r_0)+cr_1.
 \tag{3}
\]
A positive offer interaction need not imply $c>0$. If $r_0=0<r_1$, even observing the full deterministic arm outcomes and inputs cannot separate $a$ from $c$ without further restrictions.
\end{proposition}
\begin{proof}
The four means are $0,ar_0,b,ar_1+b+cr_1$, giving (3). For example take $a=1/4$, $b=1/4$, $r_0=1/4$, $r_1=3/4$ and $c=-1/16$. All values on the input square lie in $[0,1]$, yet (3) equals $5/64>0$. If $r_0=0$, observed inputs satisfy $P=PT$. Replacing $(a,c)$ by $(a+t,c-t)$ changes no observed outcome or input. Small choices of either sign preserve bounded scores and other strict parameter restrictions, showing observational equivalence with differing technological interactions.
\end{proof}

\section*{Endogenous parental effort and net welfare}
Consider a distinct fully specified mechanism. Preschool delivery is $p\in\{0,1\}$ and home visiting changes the real parental-effort cost coefficient from $k_0$ to $k_1$, where $k_0>k_1\ge0$. Raw developmental benefit is
\[
 F(p,t)=ap+(b+cp)t-\frac d2t^2,
\]
with $a\ge0$, $b>0$, $d>0$ and $b+c>0$. Parents choose $t\in[0,\bar t]$ to maximize $F(p,t)-k_ht^2/2$, with $\bar t>\max\{b,b+c\}/(d+k_1)$. There is no other unmeasured input. Readiness is $S=F/K$ at the chosen effort, where a fixed known $K$ is large enough to place these four outcomes in $[0,1]$. Assume the discounted monetary value of developmental benefit is $F$ itself. Thus $V=F-k_ht^2/2-C_{ph}$, with $C_{ph}$ the separately measured delivery-resource cost; parental time cost is not counted again in $C$.

\begin{theorem}[Solved effort response and divergent welfare signs]
The unique parental optimum is $t_{ph}=(b+cp)/(d+k_h)$. Put
\[
 A=(b+c)^2-b^2,\qquad R(k)=\frac{d+2k}{2(d+k)^2},\qquad
 \Delta C=C_{11}-C_{10}-C_{01}+C_{00}.
\]
The exact interactions are
\[
 \gamma_S=\frac A K\{R(k_1)-R(k_0)\},\qquad
 \gamma_V=\frac A2\left\{\frac1{d+k_1}-\frac1{d+k_0}\right\}-\Delta C.
\]
In this mechanism the readiness interaction has the sign of $c$, but a strictly positive readiness interaction can coexist with a strictly negative welfare interaction.
\end{theorem}
\begin{proof}
The parent's objective is a strictly concave quadratic with derivative $b+cp-(d+k_h)t$. The assumed effort bound makes the displayed stationary point interior and unique. Substitution gives $F(p,t_{ph})=ap+(b+cp)^2R(k_h)$ and $V_{ph}=ap+(b+cp)^2/[2(d+k_h)]-C_{ph}$. Alternating differences give the formulas. Since $R'(k)=-k/(d+k)^3$, $R(k_1)>R(k_0)$ whenever $k_0>k_1\ge0$. Also $A=c(2b+c)$ and $2b+c>0$, establishing the readiness sign. For $c>0$, choosing a complementary delivery cost $\Delta C$ larger than the positive benefit term makes $\gamma_V<0$ without changing readiness or parental effort.
\end{proof}

The last theorem's sign correspondence is model-specific and does not override the attendance counterexample. A credible empirical resolution needs all four randomized offers, implementation and actual-input measures, a supported exclusion argument before interpreting technological cross-effects, and five-year cost and welfare follow-up. Baseline heterogeneity can be reported with fixed strata; conditioning on realized participation would instead select on the response being studied. Neither short-run readiness nor a significant offer interaction supplies the missing long-run welfare valuation.

\section{Completion Estimate}
64\%

This is a post hoc, heuristic estimate for the full catalogue target.
The attempt identifies factorial offer interactions with finite-sample uncertainty, gives input-rank conditions and attendance counterexamples, and solves parental effort with divergent readiness and welfare signs. Target-district complementarity still requires all supported offers, credible actual-input exclusions, measured implementation and parental responses, spillover handling, and five-year welfare and resource-cost evidence.

\begin{thebibliography}{9}
\bibitem{duncan} G. J. Duncan, A. Kalil, M. Mogstad and M. Rege (2022), \emph{Investing in Early Childhood Development in Preschool and at Home}, Becker Friedman Institute Working Paper 2022-58; NBER Working Paper 29985. Section 7, especially printed pp. 91--93, motivates combined programs and measurement of parental responses. \url{https://bfi.uchicago.edu/wp-content/uploads/2022/05/BFI_WP_2022-58.pdf}; \url{https://doi.org/10.3386/w29985}.
\end{thebibliography}

\end{document}
